\section*{Resumo}

 O objetivo deste projeto é avaliar diferentes métodos de extração de características de impressões digitais voltados para a verificação biométrica de usuários. A metodologia do projeto estabelece o uso de conjuntos de dados públicos, como os das competições FVC (Fingerprint Verification Competition) e da Neurotechnology, e métricas de desempenho como FMR, FNMR, EER e acurácia balanceada. Até o momento, o trabalho desenvolveu uma revisão bibliográfica e iniciou a experimentação utilizando um modelo pré-treinado baseado na arquitetura~\textit{DeepPrint}, que extrai das impressões digitais uma representação de tamanho fixo, que pode ser usada em algoritmos de classificação baseados em distâncias. O modelo foi testado utilizando conjunto de dados FVC2004. Como trabalhos futuros, outras técnicas baseadas em Aprendizado de Máquina serão avaliadas nos demais conjuntos de dados da FVC.